\section{A sieved endpoint for actual root intervals}
\label{sw-section}

For positive integers $n,B$ and $B\le k<2B$, put
$w_k=3k+\zeta$, $T_n(k)=|P(w_k^n)|$, and let $t_k$ be the logarithm
of the largest side of its primitive positive sector normalization.
The preceding actual root-interval argument proves
\[
 t_k\ge n\log(3B),\qquad
 \log|P(w_k^d)|\le3d\log(6B+1).
\]
For real $5\le Y\le Z$, write
\[
 E_{Y,Z}(k)=\sum_{Y<p\le Z}(v_p(T_n(k))-3)_+\log p.
\]
These are actual integer roots and valuations, not a uniform reference
distribution substituted for a power orbit.

\begin{theorem}[A sieved actual interval estimate]\label{sw-estimate}
If $Z>3n$, and $\sigma(n)=\sum_{d\mid n}d$, then
\[
 \frac1B\sum_{B\le k<2B}\frac{E_{Y,Z}(k)}{t_k}
 \le \frac{10\log Y}{Y^3\log(3B)}
 +\frac{36Z\sigma(n)}{Bn\log(Z/(3n))}
 +\frac{\log n}{n\log(3B)}.
\]
The middle term may also be replaced by its minimum with $12(Z+1)/B$.
\end{theorem}
\begin{proof}
At an eligible prime of exact homogeneous rank $d\mid n$, the
actual root count is $3\varphi(d)-\mathbf1_{d=1}$, and each root
lifts uniquely. The first-depth positive cost summed over the interval
is therefore at most
\[
 \frac{3\varphi(d)B\log p}{p^3(p-1)}
       +9\varphi(d)d\log(6B+1).
\]
The first term follows by summing depth layers at least four; the
second uses their actual polynomial-height cap. At a fixed prime,
the sum of $\varphi(d)$ over eligible divisors of $n$ is at most $n$.
The already proved elementary prime sum gives the first term of the
asserted normalized estimate.

For the endpoint term, a rank-$d$ prime lies in a reduced class
$1$ or $-1$ modulo $3d$. The established Brun--Titchmarsh inequality
gives, for $X>q$ and $(a,q)=1$,
\[
 \pi(X;q,a)\le\frac{2X}{\varphi(q)\log(X/q)}.
\]
The precise source is Montgomery--Vaughan, \emph{The large sieve},
Theorem~2, printed page~121, equation~(1.10), in the
\href{https://personal.science.psu.edu/rcv4/personal/Publications/large_sieve.pdf}{author-hosted primary paper}.
Its interval form gives the displayed version by taking the left
endpoint down to zero. The bound is uniform for $X>q$.
Thus the two classes through $Z$ contain at most
$4Z/(\varphi(3d)\log(Z/(3d)))$ primes. Since $Z>3n$ and
$\varphi(3d)\ge2\varphi(d)$, their total endpoint cost is at most
\[
 36Z\log(6B+1)\sum_{d\mid n}
 \frac{d\varphi(d)}{\varphi(3d)\log(Z/(3d))}
 \le\frac{18Z\sigma(n)\log(6B+1)}{\log(Z/(3n))}.
\]
Divide by $Bn\log(3B)$ and use $6B+1\le(3B)^2$.
The actual additional lifting cost is at most $\log n$ at each root,
by $(s+u-3)_+\le(s-3)_++u$ and $\sum_pv_p(n)\log p=\log n$.
This proves the estimate. Minimizing with the previously proved
unsieved endpoint bound gives the last assertion.
\end{proof}

\begin{corollary}[A window beyond the root-interval length]\label{sw-beyond}
Let $n\to\infty$ with $\sigma(n)/n\le C$ for a fixed constant.
Set $B=n^4$, $Y=n^{1/6}$ and $Z=\lfloor B\sqrt{\log n}\rfloor$.
Then the mean in Theorem~\ref{sw-estimate} is
$O_C((\log n)^{-1/2})$, and the proportion exceeding
$(\log n)^{-1/4}$ is $O_C((\log n)^{-1/4})$.
\end{corollary}
\begin{proof}
The first and last terms are at most $5/(12\sqrt n)$ and $1/(4n)$.
Furthermore
$\log(Z/(3n))=3\log n+\tfrac12\log\log n-\log3+o(1)$.
This bounds the middle term as asserted. Apply Markov's inequality
to the nonnegative actual cost.
\end{proof}

Prime-power indices satisfy the divisor-sum hypothesis with $C=2$.
Here $Z/B\to\infty$; for prime indices the window extends beyond the
first possible top-rank progression, without asserting that it contains
a prime. Every prime above $Z$ and every exceptional root remain
uncontrolled. No estimate for a prescribed root's full signed packet
has been deduced from this average.
